\chapter{Unit 4: Product Design and Prototyping}
\label{chap:unit4_qb}

\section*{Practice Question Set (30 MCQs)}
\addcontentsline{toc}{section}{Practice Question Set (30 MCQs)}

\mcqitem{1}{What is the core definition of a prototype in Design Thinking?}{A finished product ready for mass commercial distribution}{A testable representation of an idea designed to answer specific questions and test assumptions}{A legal patent document}{A financial balance sheet}

\mcqitem{2}{What does the prototyping philosophy "Fail early, fail cheaply" mean?}{Intentionally destroying products to claim insurance}{Uncovering critical design flaws in low-cost prototypes before expensive production begins}{Giving up on projects at the first sign of trouble}{Lowering product prices below production costs}

\mcqitem{3}{Which of the following is a characteristic of a LOW-FIDELITY prototype?}{Real-time database connectivity}{Hand-drawn paper sketches, cardboard models, and rough mockups}{High-resolution animations and pixel-perfect branding}{Production-grade materials}

\mcqitem{4}{When should a design team primarily choose a Low-Fidelity prototype?}{When performing final regulatory compliance testing}{Early in the design process to explore concepts, layout, and user flows rapidly}{When measuring microsecond software server response latency}{For investor demonstrations of fully functional algorithms}

\mcqitem{5}{What is a key advantage of High-Fidelity clickable prototypes?}{They can be created in 30 seconds with a pen}{They allow realistic usability testing of detailed interactions and visual hierarchy}{They require zero software tools}{They prevent users from giving any negative feedback}

\mcqitem{6}{In a paper prototyping testing session, what role does the "Computer" play?}{Records video of the test on a laptop}{A facilitator who silently swaps paper screens and pop-ups in response to user taps}{Analyzes statistical metrics automatically}{Explains the backend database code to the participant}

\mcqitem{7}{What is a "Storyboard" in user experience design?}{A screenplay written for a movie studio}{A comic-strip style sequence of visual frames showing a user interacting with a product in real-world context}{A task board showing software development sprint tickets}{A list of customer testimonials}

\mcqitem{8}{What is the primary purpose of a Wireframe?}{Demonstrating color palettes and typography branding}{Providing a grayscale skeletal layout that establishes content hierarchy and interface structure}{Showing 3D electrical wiring diagrams of a building}{Writing automated database queries}

\mcqitem{9}{Why are wireframes traditionally created in grayscale (black, white, gray)?}{Because color printing is too expensive}{To keep focus strictly on layout, information hierarchy, and functionality without color distraction}{Because computers cannot display color during early design}{Because users dislike colorful interfaces}

\mcqitem{10}{In UI/UX wireframing conventions, what does a rectangle with an 'X' drawn across it represent?}{A delete button}{An image or media placeholder}{A broken database link}{A closed window}

\mcqitem{11}{When should Radio Buttons be used in an interface wireframe?}{When a user can select multiple options simultaneously}{When a user must choose exactly ONE option from a mutually exclusive list}{For scrolling through long pages}{For entering password text}

\mcqitem{12}{When should Checkboxes be used in a form?}{When a user may select zero, one, or MULTIPLE independent options}{When only one binary choice is allowed}{To trigger an immediate payment transfer}{To submit a form to the server}

\mcqitem{13}{In a user flow diagram, what does an Oval (Terminator) shape represent?}{A decision point}{The Start or End of a user workflow}{A database storage unit}{A clickable button}

\mcqitem{14}{In a user flow diagram, what does a Diamond shape represent?}{A user screen}{A conditional decision point with branching paths (e.g., Yes / No)}{The start of the process}{A payment gateway}

\mcqitem{15}{What is the "Happy Path" (Primary Path) in a user flow?}{A colorful design theme}{The default, error-free path where a user accomplishes their goal in the shortest direct sequence}{A path that only children use}{A marketing advertisement flow}

\mcqitem{16}{What is a Minimum Viable Product (MVP)?}{A stripped-down, broken version of a product that barely functions}{The smallest version of a product that delivers real value to early adopters and enables maximum validated learning}{A product with every single conceivable feature included}{A product built with zero budget}

\mcqitem{17}{Which of the following is a classic analogy for an MVP?}{Delivering a single wheel, then an axle, then a car chassis}{Delivering a skateboard first, then a bicycle, then a motorcycle, then a car (each stage is functional transport)}{Delivering a car without brakes}{Delivering an engine sitting on the floor}

\mcqitem{18}{What is a "Landing Page MVP"?}{An airport runway management software}{A single web page describing a value proposition with a "Sign Up" button to gauge consumer demand before building}{A fully coded e-commerce website with 50,000 products}{A printed brochure handed out on the street}

\mcqitem{19}{What is a "Concierge MVP"?}{A software product for luxury hotel receptionists}{Delivering the product service manually and personally behind the scenes to learn user needs before automating}{An automated AI chatbot}{A physical prototype made of titanium}

\mcqitem{20}{What is a "Wizard of Oz" prototype?}{A fantasy video game prototype}{A prototype that looks completely automated to the front-end user but has humans manually operating backend tasks}{A prototype that changes color dynamically}{A prototype built using voice-over actors}

\mcqitem{21}{Why is testing with representative content (realistic text/prices) superior to "Lorem Ipsum" in wireframes?}{Lorem Ipsum is illegal in software design}{Realistic content exposes actual spacing constraints, text overflow, and cognitive comprehension issues}{Realistic text makes the wireframe look like a finished website}{It prevents copyright violations}

\mcqitem{22}{Which prototype fidelity is best suited for testing complex micro-interactions, like swipe gestures and animated transitions?}{Paper sketch}{Cardboard cutout}{High-Fidelity interactive digital prototype}{Static text document}

\mcqitem{23}{In digital product design, what does an "Empty State" design represent?}{A screen when the user's internet is disconnected}{What a user sees when there is no data to display yet (e.g., "No orders placed yet"), guiding them to take action}{A completely blank white screen due to a server crash}{A logout confirmation screen}

\mcqitem{24}{What is the riskiest assumption for a startup building an on-demand laundry app?}{What shade of blue to use for the logo}{Whether customers are willing to trust strangers to wash and return their clothes reliably}{Whether to use React or Angular for the web frontend}{The font size of the copyright notice}

\mcqitem{25}{How many user testing participants are typically recommended by usability expert Jakob Nielsen for discovering ~85\% of usability errors in a prototype?}{100 users}{Exactly 5 users per target user group}{1,000 users}{Only 1 user}

\mcqitem{26}{What does "Traceability" mean in product design?}{Tracing drawings on transparent paper}{The ability to track every wireframe component and feature back to a validated user need and requirement}{Recording user GPS locations}{Following competitors' design styles}

\mcqitem{27}{In wireframing, what is the role of an "Annotation"?}{Artistic decoration in the margins}{Numbered descriptive notes explaining dynamic behaviors, validation rules, and error states}{User comments on social media}{Software pricing details}

\mcqitem{28}{Which of the following is an example of an Alternate Path in a checkout user flow?}{User enters address $\rightarrow$ enters card $\rightarrow$ payment succeeds}{User's credit card fails $\rightarrow$ system prompts for another payment method $\rightarrow$ user retries}{User searches product $\rightarrow$ views product}{User opens app $\rightarrow$ logs in}

\mcqitem{29}{What is the main danger of building a High-Fidelity prototype too early in a project?}{It uses too much electricity}{Teams become emotionally attached to the visual polish and resist making fundamental structural changes}{Users refuse to test beautiful prototypes}{High-fidelity prototypes cannot be shared online}

\mcqitem{30}{What constitutes a complete, testable MVP hypothesis?}{"We believe [target user] has [problem], and will achieve [outcome] using [feature/solution], measured by [metric]."}{"We believe our app will make millions of dollars by next month."}{"We believe everyone in the world needs our new software."}{"We believe technology is the future of business."}


\newpage
\section*{Answer Key \& Step-by-Step Explanations}
\addcontentsline{toc}{section}{Answer Key \& Explanations}

\begin{answerkeytable}{Unit 4: Quick Answer Key}
\begin{tabular}{c c c c c c c c c c}
\toprule
\textbf{Q1} & \textbf{Q2} & \textbf{Q3} & \textbf{Q4} & \textbf{Q5} & \textbf{Q6} & \textbf{Q7} & \textbf{Q8} & \textbf{Q9} & \textbf{Q10} \\
\textbf{B} & \textbf{B} & \textbf{B} & \textbf{B} & \textbf{B} & \textbf{B} & \textbf{B} & \textbf{B} & \textbf{B} & \textbf{B} \\
\midrule
\textbf{Q11} & \textbf{Q12} & \textbf{Q13} & \textbf{Q14} & \textbf{Q15} & \textbf{Q16} & \textbf{Q17} & \textbf{Q18} & \textbf{Q19} & \textbf{Q20} \\
\textbf{B} & \textbf{A} & \textbf{B} & \textbf{B} & \textbf{B} & \textbf{B} & \textbf{B} & \textbf{B} & \textbf{B} & \textbf{B} \\
\midrule
\textbf{Q21} & \textbf{Q22} & \textbf{Q23} & \textbf{Q24} & \textbf{Q25} & \textbf{Q26} & \textbf{Q27} & \textbf{Q28} & \textbf{Q29} & \textbf{Q30} \\
\textbf{B} & \textbf{C} & \textbf{B} & \textbf{B} & \textbf{B} & \textbf{B} & \textbf{B} & \textbf{B} & \textbf{B} & \textbf{A} \\
\bottomrule
\end{tabular}
\end{answerkeytable}

\begin{expbox}[Detailed Pedagogical Explanations]
\begin{enumerate}[leftmargin=6mm, itemsep=2.5mm]
  \item \textbf{[Q1: Option B]} A prototype is an exploratory learning tool to test specific assumptions.
  \item \textbf{[Q2: Option B]} Early cheap failures reveal wrong assumptions before major capital is committed.
  \item \textbf{[Q3: Option B]} Lo-Fi prototypes are quick, rough representations like paper sketches and cardboard.
  \item \textbf{[Q4: Option B]} Lo-Fi is ideal early on for exploring structures, workflows, and core concepts.
  \item \textbf{[Q5: Option B]} Hi-Fi prototypes provide realistic interaction for testing fine UI and usability details.
  \item \textbf{[Q6: Option B]} In paper prototyping, the "Computer" physically manipulates paper states based on user actions.
  \item \textbf{[Q7: Option B]} A storyboard visually illustrates the user journey and context across comic-like panels.
  \item \textbf{[Q8: Option B]} Wireframes establish structural layout, visual hierarchy, and interface functional zones.
  \item \textbf{[Q9: Option B]} Grayscale keeps feedback focused on structure and function rather than subjective color debates.
  \item \textbf{[Q10: Option B]} A crossed rectangle is the universal design blueprint placeholder for images/video.
  \item \textbf{[Q11: Option B]} Radio buttons represent single-choice mutually exclusive selections.
  \item \textbf{[Q12: Option A]} Checkboxes allow zero, one, or multiple independent selections.
  \item \textbf{[Q13: Option B]} Ovals (terminators) denote the starting and ending points of a flowchart.
  \item \textbf{[Q14: Option B]} Diamond shapes represent conditional decision logic (e.g., Is user logged in?).
  \item \textbf{[Q15: Option B]} The happy path is the primary frictionless route to task completion.
  \item \textbf{[Q16: Option B]} An MVP is the leanest functional product that delivers core value and generates learning.
  \item \textbf{[Q17: Option B]} An MVP must provide functional value at every stage (skateboard $\rightarrow$ car).
  \item \textbf{[Q18: Option B]} Landing page MVPs test market demand and click interest before writing code.
  \item \textbf{[Q19: Option B]} Concierge MVPs deliver services manually to learn direct customer needs.
  \item \textbf{[Q20: Option B]} Wizard of Oz prototypes simulate automated backends with hidden manual human labor.
  \item \textbf{[Q21: Option B]} Real content exposes actual text wrapping, button sizing, and cognitive clarity.
  \item \textbf{[Q22: Option C]} Hi-Fi digital prototypes are required to test animated transitions and responsive gestures.
  \item \textbf{[Q23: Option B]} Empty states guide first-time users when no data exists yet.
  \item \textbf{[Q24: Option B]} Customer trust in handing over laundry is the fundamental value assumption.
  \item \textbf{[Q25: Option B]} Jakob Nielsen's research showed 5 users uncover ~85\% of usability issues.
  \item \textbf{[Q26: Option B]} Traceability connects UI components directly back to user requirements.
  \item \textbf{[Q27: Option B]} Annotations document dynamic behavior, data validation, and logic for developers.
  \item \textbf{[Q28: Option B]} Payment retry/recovery represents an alternate/exception flow path.
  \item \textbf{[Q29: Option B]} High early polish causes sunk cost fallacy and attachment to flawed structures.
  \item \textbf{[Q30: Option A]} A structured hypothesis links user, problem, proposed solution, and measurable metric.
\end{enumerate}
\end{expbox}
